\documentclass{article}
\usepackage{amsmath,amsthm,amssymb}
\usepackage[margin=1in]{geometry}

\newcommand{\defc}{\operatorname{defc}}
\newcommand{\dinv}{\operatorname{dinv}}
\newcommand{\area}{\operatorname{area}}
\newcommand{\NRCM}{\operatorname{NRCM}}

\newtheorem{conjecture}{Conjecture}
\newtheorem{proposition}{Proposition}

\title{Dyck Skeleton String Decompositions}
\author{}
\date{}

\begin{document}
\maketitle

\section{Overview}

This note studies string decompositions inside fixed deficit layers.  A string
is a sequence of paths along which area increases by \(1\), dinv decreases by
\(1\), and deficit remains fixed.  The guiding example is Hawkes's
low-deficit Dyck skeleton formula, where special skeletons start strings whose
lower halves cover the corresponding deficit layers.

The main conjectural extension considered here is the case \(r=\tau s+1\).
For this family we define \(\tau\)-Dyck sequences, a \(\tau\)-analogue of
dinv, special skeletons, and a conjectural lower-half string formula.  A full
coefficient formula requires either deficit-layer \(q,t\)-symmetry or a direct
full coefficient-dictionary identity.

The final sections describe finite computational checks.  These checks support
the conjectural formula in selected ranges and also test a separate
general-slope diagnostic called NRCM.

All coefficient polynomials in this note use the convention
\(q^{\area}t^{\dinv}\).  Thus the lower half of a fixed deficit layer is the
part with \(\area\le\dinv\).

The rest of the note first introduces the mathematical sequences, paths, and
claims and then describes the finite checks attached to them.

\section{Common Vocabulary}

For a fixed total degree \(M\), the deficit of a sequence or path \(X\) is
\[
  \defc(X)=M-\area(X)-\dinv(X).
\]
A string of fixed deficit \(e\) is a sequence
\[
  X_0,X_1,\ldots,X_m
\]
such that \(\area(X_{i+1})=\area(X_i)+1\),
\(\dinv(X_{i+1})=\dinv(X_i)-1\), and \(\defc(X_i)=e\) for every \(i\).
Thus every term of the string has total degree \(M-e\).

The lower half of deficit \(e\) means the paths with
\[
  \area\le \left\lfloor\frac{M-e}{2}\right\rfloor.
\]
A lower-half decomposition is a disjoint cover of this lower-half set by the
initial segments of strings.

The remaining sections use this vocabulary in three settings: the classical
Dyck setting, the \(r=\tau s+1\) rational setting, and the NRCM diagnostic.
The mathematical definitions are stated first.  The scripts and finite checks
are described separately afterward.

\section{The Classical Model}

An ordinary Dyck sequence of length \(n\) is an integer sequence
\[
  x=(x_0,\ldots,x_{n-1})
\]
with
\[
  x_0=0,\qquad x_i\ge0,\qquad x_{i+1}\le x_i+1.
\]
Its area is
\[
  \area(x)=\sum_{i=0}^{n-1}x_i,
\]
and the classical dinv statistic in this notation is
\[
  \dinv(x)=\#\{(i,j):0\le i<j<n,\ x_i=x_j\text{ or }x_i=x_j+1\}.
\]
Here \(M=\binom n2\), so
\[
  \defc(x)=\binom n2-\area(x)-\dinv(x).
\]

The skeleton terminology in this section is imported from Section 5
of~\cite{Hawkes2026DyckSymmetric}.  In Hawkes's setup, extractable entries are
the entries removed by the recursive down move, full Dyck skeletons are the
terminal Dyck sequences with no extractable entry, and special Dyck skeletons
are the permitted string starts after Hawkes's special attachment convention.
The \(\mathrm{up}\) and \(\mathrm{down}\) moves then generate the strings.

The classical theorem used as the template here is Hawkes's low-deficit
skeleton formula.  All formulas in this note use the convention
\(q^{\area}t^{\dinv}\):
\[
  \left.C_n(q,t)\right|_{\binom n2-2n+8\le \deg_{q,t}\le \binom n2}
  =
  \sum_{\substack{S\text{ special Dyck skeleton of length }n\\
                  \defc(S)\le 2n-8}}
  \frac{q^{\area(S)}t^{\dinv(S)+1}
        -q^{\dinv(S)+1}t^{\area(S)}}{t-q}.
\]
The rational expression is a finite interval contribution: if
\(a=\area(S)\), \(\nu=\dinv(S)\), and \(e=\defc(S)\), then in the low-deficit
range of the theorem one has \(a\le\nu\), and
\[
  \frac{q^a t^{\nu+1}-q^{\nu+1}t^a}{t-q}
  =
  \sum_{i=a}^{\nu}q^i t^{M-e-i}.
\]
Hawkes proves that these strings cover the lower half of every deficit layer
with \(\defc\le2n-8\).  Equivalently, their lower-half contributions are
truncated at the midpoint:
\[
  \sum_{i=a}^{\left\lfloor(M-e)/2\right\rfloor} q^i t^{M-e-i}.
\]
The upper half is then supplied by the symmetry \(C_n(q,t)=C_n(t,q)\).

\section{The \(\tau\)-Dyck Model for \(r=\tau s+1\)}

Fix integers \(s\ge1\) and \(\tau>1\), and put \(r=\tau s+1\).  For the
conjectural model in this note, the ordinary Dyck growth bound \(+1\) is
replaced by \(+\tau\).

A \(\tau\)-Dyck sequence of length \(s\) is an integer sequence
\[
  x=(x_1,\ldots,x_s),\qquad
  x_1=0,\qquad x_i\ge0,\qquad x_{i+1}\le x_i+\tau.
\]
These inequalities imply
\[
  0\le x_i\le \tau(i-1).
\]
Thus there are only finitely many \(\tau\)-Dyck sequences of each length.

The area statistic is
\[
  \area(x)=\sum_{i=1}^s x_i.
\]
The dinv statistic used in this note is the following \(\tau\)-analogue:
\[
  \dinv(x)=\sum_{1\le i<j\le s}
  \begin{cases}
    \max(0,x_i+\tau-x_j), & x_i\le x_j,\\
    \max(0,x_j+1+\tau-x_i), & x_i>x_j.
  \end{cases}
\]
The total degree is
\[
  M=\tau\binom{s}{2},
\]
and
\[
  \defc(x)=M-\area(x)-\dinv(x).
\]

For a deficit cutoff \(B\), define the aggregate polynomial
\[
  \mathcal C_{\tau,s}^{\le B}(q,t)
  =
  \sum_{\substack{x\text{ is a }\tau\text{-Dyck sequence}\\
                  0\le \defc(x)\le B}}
      q^{\area(x)}t^{\dinv(x)}.
\]
The conjectural bound is
\[
  B=(s-2)(\tau+1)-4.
\]

For comparison with the position-coordinate notation used later in the NRCM
diagnostic, when \(r=\tau s+1\) one has
\(\lfloor ri/s\rfloor=\tau i\) for \(0\le i<s\), so
\(\Delta_i=\tau\) for \(0\le i<s-1\).  In \(Q\)-coordinates, path-validity
gives
\[
  Q_{i+1}\le Q_i+\tau\qquad(0\le i<s-1).
\]
Identifying \(x_{i+1}=Q_i\) gives exactly the \(\tau\)-Dyck growth bound
\(x_{i+2}\le x_{i+1}+\tau\).

\section{Rational Special Skeletons}

For this conjectural model, all skeletons are \(\tau\)-Dyck sequences
of the fixed length \(s\).  Deletion is used only to define extractability and
the recursive map operations.

Let \(x=(x_1,\ldots,x_s)\) be a \(\tau\)-Dyck sequence.  An entry \(x_j\) is extractable if:
\[
\begin{gathered}
  x_j>0,\\
  \#\{i<j:\max(0,x_j-\tau)\le x_i<x_j\}=1,
\end{gathered}
\]
and deleting \(x_j\) leaves a \(\tau\)-Dyck sequence.  The endpoint
conventions are explicit:
\begin{itemize}
  \item \(j=1\) cannot occur because \(x_1=0\);
  \item if \(1<j<s\), deletion is equivalent to the splice inequality
  \[
    x_{j+1}\le x_{j-1}+\tau;
  \]
  \item if \(j=s\), there is no extra splice inequality.
\end{itemize}
The sequence after deletion is
\[
  (x_1,\ldots,x_{j-1},x_{j+1},\ldots,x_s).
\]
The recursive moves use left-to-right extraction order: the first extractable
entry is the one removed.

A full skeleton is a \(\tau\)-Dyck sequence of length \(s\) with no extractable entry.
For \(s\ge4\), one full skeleton is excluded from the starting set:
\[
  (0,0,1,\underbrace{0,\ldots,0}_{s-4\text{ zeros}},\tau).
\]
For \(s=4\) this is \((0,0,1,\tau)\).  For \(s<4\), there is no excluded full
skeleton and no special upward attachment.  A special skeleton is a full
skeleton other than this excluded sequence.

The excluded full skeleton is not removed from the strings.  It is attached to
the string beginning at
\[
  (0,\ldots,0,\tau)
\]
by the special upward move
\[
  (0,\ldots,0,\tau)
  \longmapsto
  (0,0,1,\underbrace{0,\ldots,0}_{s-4\text{ zeros}},\tau)
  \qquad(s\ge4).
\]

\section{Conjectural Lower-Half Formula and Symmetry Completion}

\begin{conjecture}[Special-skeleton lower-half formula for \(r=\tau s+1\)]
Fix \(s\ge1\) and \(\tau>1\), and set \(M=\tau\binom{s}{2}\) and
\[
  B=(s-2)(\tau+1)-4.
\]
For \(0\le d\le B\), put
\[
  L_d=\left\lfloor\frac{M-d}{2}\right\rfloor.
\]
If \(B\ge0\), the lower-area part of each deficit layer \(d\le B\) is obtained
by summing over all special skeletons \(z\) of length \(s\) with
\(\defc(z)=d\).  These skeletons are lower-half starts.  Write
\[
  a=\area(z),\qquad \nu=\dinv(z),\qquad d=\defc(z).
\]
Then \(z\) contributes the lower-half interval ending at the cutoff \(L_d\):
\[
  \sum_{i=a}^{L_d} q^i t^{M-d-i}.
\]
The asserted identity is
\[
  \sum_{\substack{x\text{ is a }\tau\text{-Dyck sequence}\\
                  \defc(x)=d,\ \area(x)\le L_d}}
      q^{\area(x)}t^{\dinv(x)}
  =
  \sum_{\substack{z\text{ special skeleton}\\ \defc(z)=d}}
      \sum_{i=\area(z)}^{L_d} q^i t^{M-d-i}.
\]
\end{conjecture}

The full aggregate polynomial \(\mathcal C_{\tau,s}^{\le B}(q,t)\) follows
from this lower-half statement only with an additional symmetry input.  The
needed symmetry is deficit-layer symmetry: for each \(0\le d\le B\),
\[
  \sum_{\defc(x)=d}q^{\area(x)}t^{\dinv(x)}
  =
  \sum_{\defc(x)=d}q^{\dinv(x)}t^{\area(x)},
\]
where the sums run over \(\tau\)-Dyck sequences of length \(s\).
Equivalently, one may prove the full coefficient dictionary directly.  Without
one of these two inputs, the displayed conjecture is only a lower-half
string-start formula.

A useful algebraic packaging of a full string interval is the rational
expression
\[
  \frac{q^a t^{\nu+1}-q^{\nu+1}t^a}{t-q}.
\]
When \(a\le\nu\), this expands as the full interval
\[
  \sum_{i=a}^{\nu}q^i t^{M-d-i}.
\]
If \(a>\nu\), the same expression expands with the opposite sign over the open
interval between \(\nu\) and \(a\).  That signed expansion is algebraic
bookkeeping, not the lower-half string interpretation above.

\section{Computations for the Conjecture}

The relevant scripts are
\[
  \texttt{code/check\_r1mod\_skeleton\_strings.py}
  \qquad\text{and}\qquad
  \texttt{code/run\_official\_r1mod\_checks.py}.
\]
Some script options and output fields use the legacy name \texttt{defect}.  The
relevant names here are
\[
  \texttt{--max-defect}
  \qquad\text{and}\qquad
  \texttt{initial\_passing\_defect\_range}.
\]
In the mathematical exposition these refer to deficit.

The single-case command
\[
  \texttt{python code/check\_r1mod\_skeleton\_strings.py --tau T --s S}
\]
uses \(B=(S-2)(T+1)-4\) unless \(\texttt{--max-defect}\) is supplied.  It
performs a finite formula check by comparing two coefficient dictionaries.  The
direct dictionary is obtained by exhaustive \(\tau\)-Dyck sequence enumeration in the
retained deficit range, with pruning only when an exact suffix bound proves that
the remaining suffix cannot reach the required total \(\area+\dinv\).  The
predicted dictionary in the script is obtained from the signed rational
expressions above, so this check is a direct full-dictionary identity rather
than only a lower-half comparison.  In the exposition here, the signed terms
should be read as algebraic bookkeeping; the string picture is the lower-half
formula plus deficit-layer symmetry or direct full-dictionary agreement.  The
formula status is \(\texttt{PASS}\) iff the two dictionaries are identical.  If
the requested bound is negative, there is no nonnegative deficit layer to test.

For \(S\ge5\), the same script also attempts a lower-half map check.  The
implemented local cases are named East3, East5, West3, and West5 in the code.
This note treats those labels as implementation-defined local moves.  The
special-skeleton map check tests the implemented up and down moves:
preservation of deficit, area change by \(1\), inverse behavior on checked
steps, and disjoint coverage of the lower-half target.  A map status of
\(\texttt{PASS}\) means that the checked lower-half strings cover the target
through the supported local cases.  A status of \(\texttt{PARTIAL}\) means
that no contradiction was found before the implemented moves reached the code
label \(\texttt{unsupported\_level\_7}\).  For \(s\ge4\), this map check also
includes the special upward attachment from \((0,\ldots,0,\tau)\) to the
excluded full skeleton.

The lower-cutoff roots also have a partition-indexed count.  The accompanying
code verifies the following refined form: for \(\defc\le (s-2)\tau-1\),
special skeleton roots of deficit
\(d\) and area \(a\) are counted by partitions of size \(d\) with \(a\) parts
and with all parts at most \(s-2\).  Equivalently, after conjugating
partitions this can be phrased using partitions of length at most \(s-2\), but
then the root area records the largest part of the conjugate partition rather
than its number of parts.

The batch runner is configured to run the following finite grid.  Since no run
log is included here, this table records the intended finite test range, not
independent evidence that each listed command has been freshly run.
\[
\begin{array}{c|c|c}
\tau & \text{formula cases} & \text{formula+map cases}\\
\hline
2 & 1\le s\le14 & 5\le s\le14\\
3 & 1\le s\le12 & 5\le s\le12\\
4 & 1\le s\le10 & 5\le s\le10\\
5 & 1\le s\le9  & 5\le s\le9
\end{array}
\]
The reproducible command is
\[
  \texttt{python code/run\_official\_r1mod\_checks.py}.
\]

For example, at \((\tau,s)=(2,15)\) the conjectural bound is exactly
\[
  (15-2)(2+1)-4=35.
\]
The command
\[
  \texttt{python code/check\_r1mod\_skeleton\_strings.py --tau 2 --s 15 --report-level7}
\]
is designed to test the formula through precisely the conjectural bound
\(\defc\le35\) and reports any level-7 obstructions encountered by the partial
map.

\section{The Naive Rational Cyclic Map}

The NRCM diagnostic applies to coprime positive integers \(r\) and \(s\) with
\(s>1\).  Let
\[
  H_i=\left\lfloor\frac{ri}{s}\right\rfloor,\qquad
  L_i\in\{0,\ldots,s-1\},\quad L_i\equiv ri\pmod s,\qquad
  \Delta_i=H_{i+1}-H_i
\]
for \(0\le i<s\), where \(\Delta_i\) is used only for \(0\le i<s-1\).
A position-coordinate path is an integer sequence
\[
  Q=(Q_0,\ldots,Q_{s-1})
\]
with
\[
  Q_0=0,\qquad 0\le Q_i\le H_i.
\]
It is path-valid when the path heights
\[
  P_i=H_i-Q_i
\]
satisfy
\[
  P_0\le P_1\le\cdots\le P_{s-1}.
\]

The diagnostic area is
\[
  \area(Q)=\sum_{i=0}^{s-1}Q_i,
\]
and \(M=\sum_i H_i\).  The diagnostic deficit is defined by
\[
  \defc(Q)=\sum_{1\le i<j<s}\delta_{ij}(Q),
\]
where the summand is as follows.  Put
\[
  u_{ij}=|Q_i-Q_j|.
\]
If \(Q_i\ne Q_j\) and the comparisons \(Q_i>Q_j\) and \(L_i>L_j\) have
opposite truth values, replace \(u_{ij}\) by \(u_{ij}-1\); then replace it by
\(\max(u_{ij},0)\).  Define
\[
  v_{ij}=
  \begin{cases}
    \Delta_i-(Q_{i+1}-Q_i), & Q_i>Q_j,\\
    \Delta_{i-1}-(Q_i-Q_{i-1}), & Q_j>Q_i,\\
    0, & Q_i=Q_j.
  \end{cases}
\]
Then
\[
  \delta_{ij}(Q)=\min(u_{ij},v_{ij}).
\]
The second quantity \(v_{ij}\) is not separately clamped in this diagnostic
definition.  On path-valid inputs it is nevertheless nonnegative: the two
nonzero cases are exactly the adjacent inequalities
\(Q_{i+1}-Q_i\le\Delta_i\) and \(Q_i-Q_{i-1}\le\Delta_{i-1}\).  No replacement
by \(\max(v_{ij},0)\) is part of the implemented statistic.
Finally,
\[
  \dinv(Q)=M-\area(Q)-\defc(Q).
\]
These are the statistics used by the NRCM scripts.

For a suffix \(I_k=\{k,k+1,\ldots,s-1\}\), list its columns in increasing
label order \(L_i\).  The candidate \(T_k(Q)\) cyclically moves the \(Q\)-value
from each listed column to the next listed column and moves the last listed
value to the first listed column after adding \(1\).  Columns outside \(I_k\)
are unchanged.  The strict NRCM tries \(k=1,2,\ldots,s-1\), stops at the first
\(k\) for which \(T_k(Q)\) satisfies the capacity inequalities
\[
  0\le T_k(Q)_i\le H_i,
\]
and defines
\[
  \NRCM(Q)=T_k(Q)
\]
only if this first capacity-valid candidate is also path-valid.  If no suffix
is capacity-valid, or if the first capacity-valid suffix is not path-valid,
then \(\NRCM(Q)\) is undefined.

Here is a small successful move.  For slope \(5/3\),
\[
  H=(0,1,3),\qquad L=(0,2,1).
\]
Let \(Q=(0,1,1)\).  For \(k=1\), the suffix columns \(\{1,2\}\) appear in
label order \(2,1\).  Thus \(Q_2=1\) moves to column \(1\), and \(Q_1=1\)
wraps to column \(2\) as \(2\), giving
\[
  T_1(Q)=(0,1,2).
\]
This candidate is capacity-valid and path-valid.

The failure mode in the definition is also important.  For slope \(4/3\),
\[
  H=(0,1,2),\qquad L=(0,1,2),
\]
and \(Q=(0,0,1)\).  The first capacity-valid suffix is \(k=2\), giving
\[
  T_2(Q)=(0,0,2).
\]
Its path heights are \((0,1,0)\), which are not nondecreasing.  Strict NRCM is
therefore undefined on this \(Q\).  In this example \(k=2\) is already the
final suffix, so there is no valid NRCM value.

\section{NRCM Lower-Half Diagnostics}

The NRCM scripts are
\[
  \texttt{code/check\_nrcm\_lower\_half.py}
  \qquad\text{and}\qquad
  \texttt{code/check\_nrcm\_domain.py}.
\]
The lower-half script checks, in finite deficit layers, whether every
below-midline source has a defined NRCM image, whether the image has the same
diagnostic deficit and area one larger, whether the image remains in the
allowed target, and whether the map is injective on the checked sources.  The
domain script is narrower: it checks only definedness on below-midline
sources.

By definition, every defined strict NRCM move raises area by \(1\): the suffix
values are cyclically permuted and the wrapped value is increased by \(1\).
The proof-status issue concerns deficit preservation.  A proof that defined
strict NRCM preserves the diagnostic deficit exists outside this note and has
been partially human-verified, but it is not part of the present argument.  The
finite diagnostics below locate the deficit layers where the directed-chain
picture can be checked directly.

\begin{proposition}[Conditional NRCM chain consequence]
Fix a slope \(r/s\), a deficit \(e\), and \(M=\sum_iH_i\).  Let
\[
  \mathcal L_e=\{Q:\ Q\text{ is valid},\ \defc(Q)=e,\ 
    \area(Q)\le \lfloor(M-e)/2\rfloor\}
\]
be the checked lower-half target, including midpoint targets when
\((M-e)/2\) is an integer.  Suppose strict NRCM is defined on every valid path
of deficit \(e\) with
\[
  \area(Q)<\frac{M-e}{2},
\]
so no move is required from a midpoint target.  Suppose the checked moves
preserve deficit, raise area by \(1\), land in \(\mathcal L_e\), and are
injective on those sources.  Suppose also that every path in \(\mathcal L_e\)
which is not declared NRCM-minimal has a checked predecessor in the same
deficit layer.  Then \(\mathcal L_e\) is covered by directed NRCM chains.  The
chain starts are the NRCM-minimal paths, meaning valid paths of deficit \(e\)
that are not the NRCM image of any lower-area path in the same checked layer.
\end{proposition}

To turn such lower-half strings into a full coefficient formula, one
additionally needs equality of the predicted interval dictionary with the full
direct coefficient dictionary, or a separately proved \(q,t\)-symmetry for the
homogeneous deficit layer being considered.

The output field
\[
  \texttt{initial\_passing\_defect\_range: 0..D}
\]
means that every deficit layer \(0,1,\ldots,D\), inclusive, passed the specific
diagnostic run.  For the lower-half script this includes definedness, deficit
preservation, area increase, target membership, and injectivity.  For the
domain script it includes only definedness.  The notes accompanying these
values record tentative patterns; claims that use those patterns should include
the corresponding command and run output.

\section{Status of Claims}

\begin{center}
\begin{tabular}{p{0.31\linewidth}|p{0.54\linewidth}}
claim & status\\
\hline
classical Dyck skeleton formula &
proved by Hawkes, Theorem 5.35 of~\cite{Hawkes2026DyckSymmetric}, for
\(n\ge4\) and \(\defc\le 2n-8\)\\
\(r=\tau s+1\) special-skeleton lower-half formula &
conjectural; finite exhaustive coefficient checks are implemented for selected
\((\tau,s)\), and full recovery also needs deficit-layer symmetry or direct
full-dictionary agreement\\
East3/East5/West3/West5 lower-half map &
finite computational diagnostic using the local moves named in the code; the
implemented map stops when an unsupported level-7 case is reached\\
NRCM deficit preservation &
conditional in this note outside the finite runs; defined NRCM moves raise area
by \(1\) by construction, while preservation of diagnostic deficit is the
separate proof-status claim
\end{tabular}
\end{center}

The NRCM material is diagnostic rather than a complete theorem in this note.
The definitions, examples, and finite checks are included because they test the
same directed-chain picture for general coprime slopes.  Claims depending on
NRCM deficit preservation should be read as conditional unless accompanied by a
separate proof or by the relevant finite run output.

\begin{thebibliography}{9}

\bibitem{Hawkes2026DyckSymmetric}
Graham Hawkes,
\emph{Dyck Symmetric Functions and Applications to \(q,t\)-Catalan
Polynomials},
arXiv:2605.13003, 2026.
\[
  \texttt{https://arxiv.org/abs/2605.13003}
\]

\end{thebibliography}

\end{document}
